\documentclass[12pt,letterpaper]{article}
\usepackage[left=1.10in,right=1.10in,top=0.77in,bottom=0.80in,footskip=0.37in]{geometry}
\usepackage{lmodern,amsmath,amssymb,xcolor,hyperref}
\definecolor{ink}{HTML}{172827}
\definecolor{deepteal}{HTML}{145C60}
\definecolor{hairline}{HTML}{B1C5C3}
\definecolor{tint}{HTML}{F0F6F5}
\hypersetup{colorlinks=true,urlcolor=deepteal,linkcolor=deepteal,
  pdftitle={How many eigenvalues of a random sign matrix are real?},pdfauthor={Caden Young}}
\setlength{\parindent}{0pt}
\setlength{\parskip}{6.5pt}
\setlength{\emergencystretch}{1em}
\setlength{\footnotesep}{8pt}
\setlength{\skip\footins}{11pt plus 2pt}
\renewcommand{\footnoterule}{\kern-3pt{\color{hairline}\hrule width 0.17\linewidth height 0.5pt}\kern2.5pt}
\raggedbottom
\newcommand{\resultline}[1]{%
    \par\noindent\begingroup\setlength{\fboxsep}{8pt}%
    \colorbox{tint}{\makebox[\dimexpr\linewidth-2\fboxsep\relax][c]{%
        \color{deepteal}\large$\displaystyle #1$}}\endgroup\par}
\begin{document}

\color{ink}
{\raggedright\color{deepteal}\fontsize{26}{28}\selectfont\bfseries
How many eigenvalues of a\\
random sign matrix are real?\par}
{\small\sffamily Caden Young \enspace\textperiodcentered\enspace November 24, 2026\par}
\vspace{4pt}
{\color{hairline}\hrule height 0.6pt}
\vspace{5pt}

I am very excited to announce that I have just uploaded my preprint,
\emph{The Number of Real Eigenvalues of a Random Sign Matrix}, to the arXiv.
In it, I show that a typical large $n\times n$ sign matrix has about
$\sqrt{2n/\pi}$ real eigenvalues.

Jean Ginibre studied a similar Gaussian model \href{https://doi.org/10.1063/1.1704292}{in 1965}, where the entries are independent standard normal variables. In 1992, Edelman, Kostlan, and Shub showed that the expected number of real eigenvalues in the Gaussian model is asymptotic to $\sqrt{2n/\pi}$. Their
\href{https://escholarship.org/content/qt3hw4q0wc/qt3hw4q0wc.pdf\#page=21}{experiments with random sign matrices}
suggested the same leading term.

Terence Tao and Van Vu showed that the Gaussian prediction also holds for a wider class of matrices. For even $n$, \href{https://arxiv.org/html/1206.1893v7}{their theorem} applies when the entries are independent, real, have sufficiently light tails, and match the first four moments of a standard normal. 

In the discrete sign model, each entry is chosen independently and equally likely. Unfortunately, a fair sign matches the first three Gaussian moments, but the fourth power is $1$, whereas a standard normal has a fourth moment of $3$. Even this slight mismatch puts the simple sign model outside Tao and Vu's theorem. 

Vu conjectured that a typical sign matrix has on the order of $\sqrt{n}$ real eigenvalues, likely asymptotically matching the Gaussian model. This can be intuited by thinking of discreteness at high numbers as simulating continuity, and yielding similar properties - such as the semicircle law. In 2025, Yi Han proved a \href{https://arxiv.org/html/2502.13819v2}{lower bound
with the same leading constant as the Gaussian}.

My work gives the matching upper bound. If $R_n$ counts the real eigenvalues,
then
\resultline{
  \frac{R_n}{\sqrt n}\longrightarrow\sqrt{\frac{2}{\pi}}\qquad\text{in }L^1.
}
This proves Vu's conjecture and identifies the leading constant $\sqrt{2/\pi}$.

Feel free to contact me with any questions or comments! For more on the four-moment
theorem, I \emph{highly} recommend
\href{https://terrytao.wordpress.com/2012/06/11/random-matrices-universality-of-local-spectral-statistics-of-non-hermitian-matrices/}{Tao's post about his work with Vu}.

{\color{hairline}\hrule height 0.5pt}\smallskip
{\footnotesize\sffamily
\href{https://arxiv.org/abs/2611.27245}{Read the preprint on arXiv.}\par}
\end{document}
